
\medskip
\noindent\textbf{5.3.}
\phantomsection
\label{sga3:VIIA:5.3}

We are now interested in the left adjoint of the functor
\[
  A\longmapsto A_{\mathrm{Lie}}
\]
of \SGARefLink{sga3:VIIB:5.2:theorem}{5.2}. Let \(\mathfrak g\) be a restricted \(p\)-Lie algebra
over a ring \(R\) of characteristic \(p\), let
\(U(\mathfrak g)\) be the enveloping algebra of the Lie algebra
underlying \(\mathfrak g\) (cf.\ \textup{[BLie]}, I, \S2.1), and
let
\[
  i_{\mathfrak g}:\mathfrak g\longrightarrow U(\mathfrak g)
\]
(or simply \(i\)) be the canonical map.

Let \(A\) be a unital associative \(R\)-algebra. For every Lie
algebra homomorphism
\[
  \varphi:\mathfrak g\longrightarrow A_{\mathrm{Lie}},
\]
there is a unique homomorphism of unital \(R\)-algebras
\[
  \psi:U(\mathfrak g)\longrightarrow A
\]
such that \(\psi\circ i=\varphi\). Moreover, \(\varphi\) is a
homomorphism of restricted \(p\)-Lie algebras if and only if
\(\psi\) vanishes on all the elements
\[
  i(x)^p-i\bigl(x^{(p)}\bigr),
  \qquad x\in\mathfrak g.
\]

\noindent\textbf{Definition.}
The quotient of \(U(\mathfrak g)\) by the two-sided ideal
generated by the elements
\[
  i(x)^p-i\bigl(x^{(p)}\bigr)
\]
is denoted \(U_p^R(\mathfrak g)\), or simply
\(U_p(\mathfrak g)\). We denote by \(j_{\mathfrak g}\), or simply
\(j\), the composite
\[
  \mathfrak g
  \xrightarrow{i}
  U(\mathfrak g)
  \longrightarrow
  U_p(\mathfrak g).
\]
The algebra \(U_p(\mathfrak g)\) is called the
\emph{restricted enveloping algebra} of \(\mathfrak g\).

The preceding discussion gives the following result.

\medskip
\noindent\textbf{Proposition.}
\textit{For every unital associative \(R\)-algebra \(A\) and
every homomorphism of restricted \(p\)-Lie algebras
\[
  \varphi:\mathfrak g\longrightarrow A_{\mathrm{Lie}},
\]
there is a unique homomorphism of unital algebras
\[
  \psi:U_p(\mathfrak g)\longrightarrow A
\]
such that \(\psi\circ j=\varphi\). In other words, the functor
\(\mathfrak g\mapsto U_p(\mathfrak g)\) is left adjoint to the
forgetful functor \(A\mapsto A_{\mathrm{Lie}}\).}

\medskip
\noindent\textbf{5.3.1.}
\phantomsection
\label{sga3:VIIA:5.3.1}

With the notation of \SGARefLink{sga3:VIIA:5.3}{5.3}, put
\[
  \beta(x)=i(x)^p-i\bigl(x^{(p)}\bigr).
\]
For every \(y\in\mathfrak g\), \SGARefLink{sga3:VIIA:5.1}{formulas 5.1}~\SGARefLink{sga3:vii:paragraph:5-1:equation:i}{(i)} and \SGARefLink{sga3:VIIB:5.2:theorem}{5.2}~\SGARefLink{sga3:vii:paragraph:5-1:equation:i}{(i)}
give
\[
\begin{aligned}
  \beta(x)i(y)
  &=
  i(y)\beta(x)+[\beta(x),i(y)]\\
  &=
  i(y)\beta(x)
  +(\operatorname{ad}i(x))^p i(y)
  -i\bigl((\operatorname{ad}x)^p y\bigr)\\
  &=
  i(y)\beta(x).
\end{aligned}
\]
Thus \(\beta(x)\) belongs to the center of \(U(\mathfrak g)\);
in particular, the left ideal generated by the elements
\(\beta(x)\) is already two-sided.

On the other hand,
\[
  \beta(\lambda x)=\lambda^p\beta(x)
  \qquad(\lambda\in R),
\]
and \SGARefLink{sga3:VIIA:5.1}{formulas 5.1}~\SGARefLink{sga3:vii:paragraph:5-1:equation:ii}{(ii)} and \SGARefLink{sga3:VIIA:5.2}{5.2}~\SGARefLink{sga3:vii:paragraph:5-1:equation:ii}{(ii)} give
\[
  \beta(x+y)=\beta(x)+\beta(y)
  \qquad(x,y\in\mathfrak g).
\]
Consequently, if \((x_\alpha)\) is a family of generators of the
\(R\)-module \(\mathfrak g\), the left ideal generated by all
\(\beta(x)\) is already generated by the elements
\(\beta(x_\alpha)\).

\medskip
\noindent\textbf{5.3.2. Proposition.}\footnote{In this paragraph,
the editors changed the order, stating the result first and then
giving the detailed proof.}
\phantomsection
\label{sga3:VIIA:5.3.2}
\textit{Let \(\mathfrak g\) be an \(R\)-Lie algebra whose
underlying \(R\)-module is free with basis \((x_\alpha)\). Then
the restricted \(p\)-Lie algebra structures on \(\mathfrak g\)
correspond bijectively to the families \((y_\alpha)\) of elements
of \(\mathfrak g\) such that
\[
  \operatorname{ad}y_\alpha
  =
  (\operatorname{ad}x_\alpha)^p.
\]}

Indeed, suppose that \(\mathfrak g\) is endowed with a restricted
\(p\)-Lie algebra structure \(x\mapsto x^{(p)}\). By \SGARefLink{sga3:VIIB:5.2:theorem}{5.2}~\SGARefLink{sga3:vii:paragraph:5-1:equation:i}{(i)},
\SGARefLink{sga3:VIIB:0}{(0)}, and \SGARefLink{sga3:vii:paragraph:5-1:equation:ii}{(ii)}, the elements
\[
  y_\alpha=x_\alpha^{(p)}
\]
satisfy
\[
  \operatorname{ad}y_\alpha
  =
  (\operatorname{ad}x_\alpha)^p,
\]
and they determine the restricted \(p\)-Lie algebra structure.

We prove the converse. Since \(\mathfrak g\) is a free
\(R\)-module, the canonical map
\[
  i:\mathfrak g\longrightarrow U(\mathfrak g)
\]
is injective by the Poincar\'e--Birkhoff--Witt theorem
(cf.\ \textup{[BLie]}, I, \S2.7). We may therefore identify
\(\mathfrak g\) with an \(R\)-submodule of \(U(\mathfrak g)\).
Suppose that \((y_\alpha)\) is a family of elements of
\(\mathfrak g\) such that
\[
  \operatorname{ad}y_\alpha
  =
  (\operatorname{ad}x_\alpha)^p.
\]
Let
\[
  \pi:R\longrightarrow R,
  \qquad
  r\longmapsto r^p,
\]
and let \(\mathfrak g\otimes_\pi R\) be the \(R\)-Lie algebra
obtained by extension of scalars along \(\pi\).

There is an \(R\)-linear map
\[
  \gamma:\mathfrak g\otimes_\pi R
  \longrightarrow
  U(\mathfrak g)
\]
which sends \(x_\alpha\otimes_\pi1\) to
\(x_\alpha^p-y_\alpha\). Moreover, for every
\(x\in\mathfrak g\),
\[
  (\operatorname{ad}x_\alpha^p)(x)
  =
  (\operatorname{ad}x_\alpha)^p(x)
  =
  (\operatorname{ad}y_\alpha)(x).
\]
It follows that \(\gamma\) maps
\(\mathfrak g\otimes_\pi R\) into the center of
\(U(\mathfrak g)\). For \(x\in\mathfrak g\), put
\[
  x^{(p)}
  =
  x^p-\gamma(x\otimes_\pi1).
\]

Then \(x_\alpha^{(p)}=y_\alpha\) for every \(\alpha\). If
\[
  x=\sum_\alpha\lambda_\alpha x_\alpha,
\]
\SGARefLink{sga3:VIIA:5.1}{formula 5.1}~\SGARefLink{sga3:vii:paragraph:5-1:equation:ii}{(ii)}, applied inductively to the number of indices
\(\alpha\) for which \(\lambda_\alpha\ne0\), gives
\[
  x^p-\sum_\alpha\lambda_\alpha^p x_\alpha^p
  \in\mathfrak g.
\]
Denote this element by \(z\). Then
\[
  x^{(p)}
  =
  \sum_\alpha\lambda_\alpha^p y_\alpha+z,
\]
and hence \(x^{(p)}\in\mathfrak g\).

It is clear that
\[
  (\lambda x)^{(p)}
  =
  \lambda^p x^{(p)}.
\]
Since \(\gamma(x\otimes_\pi1)\) is central, we have
\(\operatorname{ad}x^{(p)}=\operatorname{ad}x^p\). Jacobson's
first formula, \SGARefLink{sga3:VIIA:5.1}{5.1}~\SGARefLink{sga3:vii:paragraph:5-1:equation:i}{(i)}, therefore yields
\[
  \operatorname{ad}x^{(p)}
  =
  (\operatorname{ad}x)^p.
\]
Finally, Jacobson's second formula, \SGARefLink{sga3:VIIA:5.1}{5.1}~\SGARefLink{sga3:vii:paragraph:5-1:equation:ii}{(ii)}, shows that the map
\(x\mapsto x^{(p)}\) satisfies condition \SGARefLink{sga3:vii:paragraph:5-1:equation:ii}{(ii)} of \SGARefLink{sga3:VIIB:5.2:theorem}{5.2}. It thus
makes \(\mathfrak g\) into a restricted \(p\)-Lie algebra. This
proves the proposition.

\medskip
\noindent\textbf{5.3.3. Proposition.}
\phantomsection
\label{sga3:VIIA:5.3.3}
\textit{Let \(\mathfrak g\) be a restricted \(p\)-Lie algebra
over \(R\) whose underlying module is free with basis
\((x_\alpha)\). Then
\[
  j:\mathfrak g\longrightarrow U_p(\mathfrak g)
\]
is injective. If \(z_\alpha=j(x_\alpha)\), then
\(U_p(\mathfrak g)\) has as a basis the monomials
\[
  \prod_\alpha z_\alpha^{n_\alpha},
  \qquad
  0\le n_\alpha<p.
\]
Here the \(n_\alpha\) are zero except for finitely many
\(\alpha\); the basis is assumed to be totally ordered, and the
products are taken in increasing order.}

Identify \(\mathfrak g\) with a submodule of the enveloping
algebra \(U(\mathfrak g)\) by means of the canonical map \(i\).
For every family
\[
  n=(n_\alpha)
\]
of nonnegative integers, all but finitely many zero, put
\[
  |n|=\sum_\alpha n_\alpha,
  \qquad
  x^n=\prod_\alpha x_\alpha^{n_\alpha}.
\]
Write
\[
  n_\alpha=m_\alpha+p\ell_\alpha,
  \qquad
  0\le m_\alpha<p,
\]
and put
\[
  T_n
  =
  \prod_\alpha
  x_\alpha^{m_\alpha}
  \beta(x_\alpha)^{\ell_\alpha},
\]
where
\[
  \beta(x)=x^p-x^{(p)}
\]
is the map \(\mathfrak g\to U(\mathfrak g)\) defined in \SGARefLink{sga3:VIIA:5.3.1}{5.3.1}.

For \(r\in\mathbb N\), let \(U^r\) be the \(R\)-submodule of
\(U(\mathfrak g)\) generated by the \(x^n\) for which
\(|n|\le r\). Since the graded ring
\[
  \bigoplus_r U^r/U^{r-1}
\]
is commutative (cf.\ \textup{[BLie]}, I, \S2.6), for every
\(n\) we have
\[
  T_n-\prod_\alpha x_\alpha^{n_\alpha}
  \in U^{|n|-1}.
\]
For every \(s\in\mathbb N\), the \(x^n\) with \(|n|=s\) form a
basis of \(U^s/U^{s-1}\) by the Poincar\'e--Birkhoff--Witt
theorem (\emph{loc.\ cit.}, \S2.7); the same is therefore true
of the \(T_n\) with \(|n|=s\).

As \(s=|n|\) varies, the \(T_n\) consequently form a basis of
\(U(\mathfrak g)\). The kernel \(J\) of the canonical map
\[
  U(\mathfrak g)\longrightarrow U_p(\mathfrak g)
\]
is the left ideal of \(U(\mathfrak g)\) generated by the central
elements \(\beta(x_\alpha)\), by \SGARefLink{sga3:VIIA:5.3.1}{5.3.1}. The \(T_n\) for which
\(\ell=(\ell_\alpha)\ne0\) therefore form a basis of \(J\), while
the \(T_n\) for which \(n_\alpha<p\) for every \(\alpha\) form a
basis of
\[
  U_p(\mathfrak g)=U(\mathfrak g)/J.
\]

\medskip
\noindent\textbf{5.3.3 bis.}
\phantomsection
\label{sga3:VIIA:5.3.3bis}

Let \(\mathfrak g\) be a restricted \(p\)-Lie algebra over \(R\),
and let \(f:R\to R'\) be an extension of the base ring. We claim
that the \(R'\)-module
\[
  R'\otimes_R\mathfrak g
\]
admits one and only one restricted \(p\)-Lie algebra structure
such that
\[
  \tag{*}\label{sga3:vii:paragraph:5-3-3-bis:equation:star}
  [\lambda\otimes x,\mu\otimes y]
  =
  \lambda\mu\otimes[x,y],
  \qquad
  (\lambda\otimes x)^{(p)}
  =
  \lambda^p\otimes x^{(p)}.
\]
In particular, it will follow that the functor
\[
  \mathfrak g\longmapsto R'\otimes_R\mathfrak g
\]
is left adjoint to restriction of scalars from \(R'\) to \(R\).

The uniqueness of the structure defined by \((*)\) is clear; we
prove its existence. If \(\mathfrak g\) is free with basis
\((x_\alpha)\), \SGARefLink{sga3:VIIA:5.3.2}{Proposition 5.3.2} gives a unique restricted
\(p\)-Lie algebra structure on the Lie algebra
\(R'\otimes_R\mathfrak g\) such that
\[
  (1\otimes x_\alpha)^{(p)}
  =
  1\otimes x_\alpha^{(p)}.
\]
This is the structure sought.

Suppose now that \(\mathfrak g\) is arbitrary. There is a
restricted \(p\)-Lie algebra \(L_0\), free as an \(R\)-module,
together with a surjective homomorphism
\[
  q_0:L_0\longrightarrow\mathfrak g.
\]
For example, we may take
\[
  L_0=R\otimes_{\mathbb F_p}\mathfrak g
\]
and let \(q_0\) be the homomorphism
\[
  \lambda\otimes x\longmapsto\lambda x;
\]
here \(\mathfrak g\) is free over the prime field
\(\mathbb F_p\). The kernel of \(q_0\) is a \(p\)-ideal of
\(L_0\), that is, an ideal of the Lie algebra \(L_0\) stable
under \(x\mapsto x^{(p)}\). Hence there is also a restricted
\(p\)-Lie algebra \(L_1\), free as an \(R\)-module, and a
homomorphism \(q_1:L_1\to L_0\) whose image is
\(\operatorname{Ker}q_0\). This gives the exact sequence
\begin{center}
  \phantomsection
  \label{sga3:VIIA:5.3.3bis-exact-sequence}
  \(\displaystyle
L_1 \xrightarrow{q_1} L_0 \xrightarrow{q_0} \mathfrak g \longrightarrow 0 .
\)

\end{center}
and therefore an exact sequence of \(R'\)-Lie algebras
\begin{center}
  \phantomsection
  \label{sga3:VIIA:5.3.3bis-basechange-exact-sequence}
  \(\displaystyle
R'\otimes_R L_1
  \xrightarrow{R'\otimes_R q_1}
R'\otimes_R L_0
  \xrightarrow{R'\otimes_R q_0}
R'\otimes_R\mathfrak g
  \longrightarrow 0 .
\)

\end{center}

Since \(R'\otimes_Rq_1\) is manifestly a homomorphism of
restricted \(p\)-Lie algebras, the kernel of
\(R'\otimes_Rq_0\) is a \(p\)-ideal. Thus the symbolic \(p\)-th
power operation on \(R'\otimes_RL_0\) induces, by passage to the
quotient, a map from \(R'\otimes_R\mathfrak g\) to itself (use
formula \SGARefLink{sga3:vii:paragraph:5-1:equation:ii}{(ii)} of \SGARefLink{sga3:VIIB:5.2:theorem}{5.2}). This map equips \(R'\otimes_R\mathfrak g\)
with the required restricted \(p\)-Lie algebra structure.

\medskip
\noindent\textbf{5.3.4.}
\phantomsection
\label{sga3:VIIA:5.3.4}

The canonical map
\[
  j_{\mathfrak g}:\mathfrak g\longrightarrow U_p(\mathfrak g)
\]
induces, for every extension \(R\to R'\) of the base ring, a
homomorphism
\[
  R'\otimes_Rj_{\mathfrak g}:
  R'\otimes_R\mathfrak g
  \longrightarrow
  R'\otimes_RU_p(\mathfrak g).
\]
It therefore induces a homomorphism
\[
  h:
  U_p(R'\otimes_R\mathfrak g)
  \longrightarrow
  R'\otimes_RU_p(\mathfrak g)
\]
such that
\[
  h\circ j_{R'\otimes_R\mathfrak g}
  =
  R'\otimes_Rj_{\mathfrak g}.
\]
The universal properties of \(R'\otimes_R\mathfrak g\) and of
the restricted enveloping algebra show immediately that \(h\) is
an isomorphism. We may therefore identify
\[
  U_p(R'\otimes_R\mathfrak g)
  \simeq
  R'\otimes_RU_p(\mathfrak g).
\]

In particular, let \(r\in R\), and let \(R'=R_r\) be the
localized ring. Then
\[
  \mathfrak g_r=R_r\otimes_R\mathfrak g
\]
has a canonical restricted \(p\)-Lie algebra structure over
\(R_r\). Consequently, the sheaf \(\widetilde{\mathfrak g}\) on
\(\operatorname{Spec}R\) is a quasi-coherent restricted
\(p\)-Lie algebra on \(\operatorname{Spec}R\). Moreover,
\[
  U_p^{R_r}(\mathfrak g_r)
  \simeq
  U_p^R(\mathfrak g)_r,
\]
so the sheaf associated with the presheaf
\[
  V\longmapsto
  U_p\bigl(\Gamma(V,\widetilde{\mathfrak g})\bigr)
\]
is quasi-coherent.

\noindent\textbf{Definition.}
More generally, let \(S\) be a scheme of characteristic \(p\),
and let \(\mathcal G\) be a quasi-coherent restricted \(p\)-Lie
algebra over \(\mathcal O_S\). The sheaf associated with the
presheaf
\[
  V\longmapsto U_p\bigl(\Gamma(V,\mathcal G)\bigr)
\]
is quasi-coherent. It is denoted
\(\mathcal U_p(\mathcal G)\) and is called the
\emph{restricted enveloping algebra} of \(\mathcal G\). If \(V\)
is affine, then
\[
  U_p\bigl(\Gamma(V,\mathcal G)\bigr)
  \simeq
  \Gamma\bigl(V,\mathcal U_p(\mathcal G)\bigr).
\]

